\begin{abstract}
Force-aware wiping requires contact acquisition, force regulation, transient
control, and surface progression. We introduce regime-adaptive risk-aware
model-predictive path integral (MPPI) planning for direct temporal-difference
model-predictive control (TD-MPC2) wiping. The method retains requested force
in latent dynamics and adapts planning costs and computation from causal
contact regimes and learned next-force predictions. Across 270 paired
factor-separated evaluations, it raises compound pass from 85/135 to 100/135
(+11.11 percentage points; block- and seed-aware 95\% interval [3.70, 18.52]),
reduces mean relative error by 1.83 points, and uses 8.0 fewer samples and
25.7~ms less median planning time. Task completion remains 132/135 versus
133/135, and both planners record two samples above 15~N. A 375-run one-factor
panel identifies force bias, 50-ms delay, and registered-normal error as
distinct sensitivities. In a zero-shot MuJoCo stress test, all 30 flat and
inclined evaluations retain the compound outcome, whereas none of 15 weakly
curved evaluations does. Regime adaptation therefore improves the
tracking--computation trade-off in simulation, while curved contact remains
the principal portability boundary.
% Insert the stable anonymous artifact URL here after double-blind deposition.
\end{abstract}
